\documentclass[journal=jacsat,manuscript=article]{achemso}

\usepackage[T1]{fontenc}
\usepackage[utf8]{inputenc}
\usepackage{amsmath}
\usepackage{amssymb}
\usepackage{booktabs}
\usepackage{multirow}
\usepackage{array}
\usepackage{graphicx}
\usepackage{xcolor}
\usepackage{hyperref}
\usepackage{siunitx}

%% Author information
\author{Ku Kang}
\affiliation[CBDRI]{CBRN Defense Research Institute, Ministry of National
Defense, Seoul 06796, Republic of Korea}

\author{Jeongyun Kim}
\affiliation[SNUCBE]{Department of Chemical and Biological Engineering,
  Seoul National University, Seoul 08826, Republic of Korea}

\author{Jin Yoo}
\affiliation[CBDRI]{CBRN Defense Research Institute, Ministry of National
Defense, Seoul 06796, Republic of Korea}

\author{Dongyoul Lee}
\affiliation{Department of Physics and Chemistry, Korea Military
Academy, Seoul 01805, Republic of Korea}

\email{dlee32@kma.ac.kr}


\title{Physics-Informed Machine Learning for Thermodynamic Stability Constants of
  Theranostic Radiometal--Chelator Complexes: Multi-Temperature Prediction
  via Physics-Constrained Van't~Hoff Neural Networks}

\abbreviations{GBM, Gradient Boosting Machine; RF, Random Forest; XGBoost, Extreme Gradient Boosting;
  QSPR, Quantitative Structure--Property Relationship;
  LOMO-CV, Leave-One-Metal-Out Cross-Validation;
  SHAP, SHapley Additive exPlanations;
  VH-NN, Van't Hoff Neural Network; DFT, Density Functional Theory;
  EDTA, Ethylenediaminetetraacetic acid;
  DOTA, 1,4,7,10-Tetraazacyclododecane-1,4,7,10-tetraacetic acid;
  NOTA, 1,4,7-Triazacyclononane-1,4,7-triacetic acid;
  DTPA, Diethylenetriaminepentaacetic acid;
  macropa, 2,2'-(3,12-bis(carboxymethyl)-6,9-dioxa-3,12-diazatetradecane-1,14-diyl)dibenzoic acid;
  HSAB, Hard and Soft Acids and Bases;
  ITC, Isothermal Titration Calorimetry;
  PRRT, Peptide Receptor Radionuclide Therapy;
  TRT, Targeted Radionuclide Therapy;
  LFER, Linear Free-Energy Relationship}

\keywords{radiometal chelation, thermodynamic stability, machine learning,
  van't~Hoff, multi-temperature prediction, SHAP, inverse design,
  DOTA, NOTA, radiopharmaceuticals}

\begin{document}

\begin{abstract}
Thermodynamic stability constants ($\log K_1$) at physiological ionic strength
govern the in vivo safety and transchelation resistance of radiometal--chelator
complexes used in targeted radionuclide therapy (TRT) and PET theranostics.
Experimental measurement across the full metal--chelator space is experimentally
prohibitive, motivating data-driven prediction.
We present a physics-informed machine learning pipeline trained on a curated dataset
of 140 conditional stability constants ($T = 25\,^\circ$C, $I = 0.1$~M;
14 radiometals, 14 chelators), augmented with 60 multi-temperature entries derived
from van't~Hoff integration of isothermal titration calorimetry (ITC) $\Delta H^\circ$
data for 12 metal--chelator pairs.
A five-model ablation study (Ridge through physics-constrained VH-NN) identifies
the Gradient Boosting Machine augmented with DFT interaction energies (A2~GBM+DFT)
as the optimal predictor ($R^2 = 0.831 \pm 0.169$; RMSE $= 2.35 \pm 0.93$ log~units;
95\%~CI [0.816, 0.878]; seeds $= 7$).
Leave-One-Metal-Out cross-validation (LOMO-CV) establishes a three-tier
extrapolation framework: interpolation ($R^2 \geq 0.6$; La, Ac, Y),
partial ($0 \leq R^2 < 0.6$; Pb), and out-of-domain ($R^2 < 0$;
In, Lu, Sc, Bi, Ga, Cu, Ba).
SHAP TreeExplainer analysis confirms log$K_1$ (EDTA reference) as the dominant
descriptor (importance $= 0.396$), followed by size mismatch (0.149) and
cavity volume (0.131), quantitatively validating the hard--soft acid--base (HSAB)
thermodynamic framework.
A physics-constrained VH-NN embedding the van't~Hoff equation as a differentiable
layer achieves $R^2 = 0.812$ for multi-temperature prediction with thermodynamic
self-consistency guaranteed by construction.
Inverse design via SHAP-guided feature optimization identifies a hydroxamate-macrocyclic
hybrid scaffold (predicted $\log K_1 = 26.0$ [24.0, 28.0]) as a priority candidate
for experimental validation in next-generation $^{177}$Lu theranostics;
synthesis and stability characterization of this scaffold are currently in preparation.
The curated dataset, analysis pipeline, and all model weights are provided in
Supporting Information to enable community-driven extension to additional
metal--chelator systems.
\end{abstract}

%%============================================================
\section{Introduction}
%%============================================================

Targeted radionuclide therapy and theranostic PET imaging have entered a transformative
clinical era, driven by the regulatory approvals of $^{177}$Lu-DOTATATE for
gastroenteropancreatic neuroendocrine tumors (NETTER-1 trial),\cite{Strosberg2017}
$^{177}$Lu-PSMA-617 for metastatic castration-resistant prostate cancer (VISION trial),\cite{Sartor2021}
and ongoing late-phase trials of $^{225}$Ac-PSMA-617\cite{Kratochwil2016} and
$^{225}$Ac-DOTATATE.\cite{Ballal2020}
These advances have catalyzed a broader theranostic roadmap encompassing
$^{64/67}$Cu, $^{68}$Ga, $^{212/213}$Bi, $^{212}$Pb, and $^{89}$Zr across
diverse solid tumor indications.\cite{Bodei2022,Herrmann2020,Morgan2023}

The clinical efficacy and safety of these agents depend critically on the thermodynamic
stability of the radiometal--chelator complex, quantified by the conditional stability
constant $\log K_1$ at physiological ionic strength ($I = 0.1$~M).
A thermodynamically labile complex undergoes transchelation to endogenous metal-binding
proteins (transferrin, albumin) or free phosphate in serum, releasing toxic free
radiometal to off-target organs.\cite{Price2014,Wadas2010,Ramogida2013}
The chelator therefore functions as a critical molecular safety gate, and
quantitative prediction of $\log K_1$ across the full metal--chelator space is a
central challenge in radiopharmaceutical design.\cite{Boros2019,Hu2022}

Comprehensive compilations of $\log K_1$ values exist in the NIST Critical Stability
Constants database\cite{MartellSmith2004} and primary literature for canonical
chelators (DOTA,\cite{deLeonRodriguez2008} NOTA,\cite{Notni2010}
DTPA,\cite{Pniok2014} EDTA, macropa\cite{Blei2023,RocaSabio2009}), but the
combinatorial metal--chelator space remains sparsely sampled.
A single full potentiometric characterization requires gram-scale ligand synthesis,
2--4 weeks of pH-metric titration, and expert analysis of competing equilibria ---
a resource cost of approximately \$15,000--\$30,000 per metal--chelator pair when
personnel and instrument time are included.\cite{Leavitt2001}
Systematic experimental coverage of a $14 \times 14$ matrix therefore demands
$\sim 196$ independent campaigns; at this throughput, comprehensive chelator
screening for a single new theranostic radionuclide would require decades of
coordinated laboratory effort.
A validated predictive model that evaluates the full combinatorial matrix in
$< 1$~CPU-second can compress this discovery timeline from years to weeks by
pre-selecting the top 5--10 high-confidence candidates for experimental validation,
representing a $> 95\%$ reduction in screening burden and enabling timely
response to emerging clinical radionuclide approvals.

Machine learning (ML) has emerged as a powerful surrogate for $\log K_1$ prediction.
Chaube et al.\ applied random forests to lanthanide--ligand affinity prediction
(Sci.\ Rep., 2020).\cite{Chaube2020}
Kanahashi et al.\ demonstrated gradient boosting for overall stability constants across
multiple metal classes.\cite{Kanahashi2022}
Zahariev et al.\ predicted $\log K_1$ by ML for targeted metal--ion selectivity
design,\cite{Zahariev2024} and Karunaratne et al.\ extended this to a comprehensive
multi-metal dataset.\cite{Karunaratne2025}
DFT-derived descriptors for radiometal--chelator interaction energies were introduced
by \"{O}zt\"{u}rk et al.\cite{Ozturk2024} as complementary physical descriptors.
Kim et al.\ recently reported ab initio-guided rational chelator selection for
theranostic radionuclides.\cite{Kim2026molpharm}
Despite this progress, three limitations persist across existing models:
(i)~purely data-driven descriptors ignore coordination chemistry first principles,
sacrificing interpretability;\cite{Kanahashi2022,Zahariev2024}
(ii)~training sets are limited ($n < 80$), producing high cross-validation variance;
and (iii)~all published models operate at a single temperature (25\,$^\circ$C),
whereas radiolabeling is performed at 65--95\,$^\circ$C and in vivo stability
must be assessed at 37\,$^\circ$C.\cite{Boros2019,Morgan2023}

Here we address all three limitations with an integrated physics-informed pipeline,
making four contributions that are distinct from all prior work in this
area:\cite{Karunaratne2025,Ozturk2024,Chaube2020,Kanahashi2022}
(C1)~a curated 140-entry $\log K_1$ dataset expanded with 60 multi-temperature
entries from van't~Hoff/ITC integration (25--65\,$^\circ$C, 12 metal--chelator pairs)
--- the first radiometal $\log K_1$ dataset with explicit temperature dimensionality;
(C2)~11 physics-informed descriptors grounded in HSAB theory,\cite{Pearson1963,Parr1983}
coordination geometry,\cite{Shannon1976,Hancock1989} and DFT interaction
energetics\cite{Becke1993,LeeYangParr1988,Weigend2005,Grimme2010,Barone1998,Neese2020}
for 22 structurally characterized pairs, providing electronic-structure-level
information absent from prior fingerprint-based QSPR approaches;
(C3)~a physics-constrained VH-NN that embeds the van't~Hoff equation as a
non-trainable differentiable layer, guaranteeing thermodynamic self-consistency at all
temperatures by construction and outputting interpretable $(\Delta H^\circ, \Delta S^\circ)$
pairs validated against ITC calorimetry --- a capability unavailable in any
existing radiometal chelation model; and
(C4)~a LOMO-CV three-tier extrapolation framework\cite{Lookman2019} that
explicitly maps interpolation vs.\ out-of-domain uncertainty for each metal,
translating model confidence into prioritized experimental resource allocation,
complemented by SHAP-guided inverse design\cite{Lundberg2017,Lundberg2020}
yielding experimentally falsifiable scaffold predictions with 95\% prediction intervals.

%%============================================================
\section{Theoretical Background}
%%============================================================

\subsection{Van't~Hoff Integration and Multi-Temperature Dataset Construction}
\label{sec:vanthoff_theory}

The temperature dependence of $\log K_1$ follows directly from integrating the
van't~Hoff equation under the assumption of temperature-independent $\Delta H^\circ$,
valid over the modest range 25--65\,$^\circ$C:\cite{Leavitt2001}
\begin{equation}
  \log K_1(T) \;=\; \log K_1(T_0)
    \;+\; \frac{\Delta H^\circ}{R \ln 10}
    \left(\frac{1}{T_0} - \frac{1}{T}\right)
  \label{eq:vant_hoff}
\end{equation}
where $R = 8.314$~J~mol$^{-1}$~K$^{-1}$ and $T_0 = 298.15$~K.
ITC provides $\Delta H^\circ$ and $\Delta S^\circ$ simultaneously with $\log K_1$
at a single temperature, enabling rigorous multi-temperature dataset augmentation.\cite{Leavitt2001}
The constant-$\Delta H^\circ$ assumption introduces an error below
$\pm 0.1$ log units over 40\,$^\circ$C for the enthalpically dominated
(exothermic) radiometal--polyaminocarboxylate systems studied here,
as confirmed by heat capacity measurements in the literature.\cite{Blei2023,Notni2010}

\subsection{Physics-Constrained VH-NN Architecture}

Standard regression networks fitting $\log K_1$ directly cannot extrapolate
thermodynamically across temperatures without additional physical constraints.
We introduce a VH-NN that outputs $(\Delta H^\circ, \Delta S^\circ)$ as latent
thermodynamic intermediates, with $\log K_1(T)$ computed by a fixed, non-trainable
differentiable layer implementing:
\begin{equation}
  \log K_1(T) \;=\; -\frac{\Delta G^\circ(T)}{RT \ln 10}
    \;=\; -\frac{\Delta H^\circ - T\Delta S^\circ}{RT \ln 10}
  \label{eq:vhnn}
\end{equation}
By propagating gradients through Eq.~\ref{eq:vhnn} during backpropagation, the
network learns physically meaningful $(\Delta H^\circ, \Delta S^\circ)$ pairs
and is thermodynamically self-consistent at all temperatures by construction.
This architecture is conceptually analogous to physics-informed neural networks
(PINNs) that embed differential equations as hard constraints.\cite{Karunaratne2025}

%%============================================================
\section{Methods}
%%============================================================

\subsection{Dataset Curation and Quality Control}

The primary dataset comprises 140 conditional stability constants $\log K_1$
($T = 25\,^\circ$C, $I = 0.1$~M NaClO$_4$ or NaNO$_3$) assembled from the
NIST Critical Stability Constants Database\cite{MartellSmith2004} and primary
literature.
The dataset spans 14 metals:
Sc$^{3+}$,\cite{Pniok2014}
Bi$^{3+}$,\cite{Milne2015}
Ga$^{3+}$,\cite{Notni2010,Simecek2012}
In$^{3+}$,
Zr$^{4+}$,\cite{Deri2015,Boros2014}
Lu$^{3+}$,\cite{deLeonRodriguez2008}
Y$^{3+}$,
Ac$^{3+}$,\cite{Thiele2017}
Pb$^{2+}$,\cite{FerrandoCliment2022}
La$^{3+}$,\cite{RocaSabio2009}
Cu$^{2+}$,
Ba$^{2+}$/Ra$^{2+}$,\cite{Blei2023}
and K$^+$;
and 14 chelators: DOTA, NOTA, DTPA, EDTA, macropa,\cite{Blei2023,RocaSabio2009}
CB-TE2A, AAZTA,\cite{Baranyai2007}
HOPO/DFO*,\cite{Boros2014,Deri2015}
PCTA,\cite{Paurova2016} TRAP,\cite{Simecek2012}
H4octapa, H4neunpa, NODA-GA, and THP.
Of the 140 entries, 57 are directly measured by potentiometry, spectrophotometry,
or ITC, and 83 are estimated via linear free-energy relationships (LFER)
anchored to the NIST database.\cite{MartellSmith2004,Hancock1989}

The multi-temperature subset (60 entries) was generated from Eq.~\ref{eq:vant_hoff}
using ITC-derived $\Delta H^\circ$ and $\Delta S^\circ$ values for 12 metal--chelator
pairs (Table~S3)\cite{Leavitt2001} at
$T \in \{25, 35, 45, 55, 65\}\,^\circ$C.
All $\log K_1$ values reported in this study are conditional constants at $I = 0.1$~M,
converted from literature values at other ionic strengths using the Davies
equation where necessary.
To prevent data leakage into the 25\,$^\circ$C cross-validation benchmarks,
the 60 multi-temperature entries were strictly withheld from the 5-fold
cross-validation used to evaluate models A0--A3; they were employed exclusively
for training and evaluating the VH-NN (A4) in the multi-temperature prediction
task described in Section~\ref{sec:multiT}, ensuring that temperature-interpolated
data do not inflate reported 25\,$^\circ$C cross-validation metrics.

\subsection{Physics-Informed Feature Engineering}

Eleven physics-based descriptors were computed for each metal--chelator pair (Table~S1):

\textbf{Metal-centered descriptors.}
Shannon effective ionic radii at the appropriate coordination number\cite{Shannon1976}
reflect the geometric demand of the metal ion.
Formal oxidation state (charge) governs electrostatic interaction energy.
The absolute chemical hardness $\eta$\cite{Parr1983} parameterizes HSAB character:
$\eta = (\text{IP} - \text{EA})/2$, where IP and EA are ionization potential and
electron affinity, respectively.\cite{Pearson1963,Parr1983,Hancock1989}
Preferred aqueous coordination number determines the number of available binding sites.

\textbf{Chelator-centered descriptors.}
Donor atom counts for oxygen (carboxylate/hydroxamate) and nitrogen (amine/imine)
encode the Lewis-base character of the ligand.
Ring size (macrocyclic cavity diameter) and a binary macrocyclic indicator quantify
the preorganization thermodynamic advantage.\cite{RocaSabio2009,Baranyai2007}

\textbf{Metal--chelator descriptors.}
Size mismatch ($\Delta r = r_\mathrm{ion} - r_\mathrm{cavity}$, \AA) captures the
geometric complementarity between the metal and the chelator cavity.
Estimated cavity volume (\AA$^3$) is computed from ring size.
The log$K_1$ (EDTA reference) exploits the established LFER: stability constants
correlate linearly with a benchmark chelator across the metal
series,\cite{Hancock1989,MartellSmith2004} providing a thermodynamic anchor.

\textbf{DFT interaction energy.}
For the 22 metal--chelator pairs with available crystal structures, gas-phase interaction
energies ($E_\mathrm{int}$, kJ~mol$^{-1}$) were computed at the
B3LYP-D3(BJ)\cite{Becke1993,LeeYangParr1988,Grimme2010}/def2-TZVP\cite{Weigend2005}
level of theory in implicit water (CPCM solvent model)\cite{Barone1998}
using ORCA 5.0.\cite{Neese2020}
Counterpoise correction (BSSE) was applied using the Boys--Bernardi scheme.
$E_\mathrm{int}$ captures short-range orbital and dispersion interactions not
encoded by classical force-field descriptors.\cite{Ozturk2024}
The DFT interaction energy descriptor is independently validated by a
companion study from our group employing the same ORCA computational framework
and CPCM(water) solvation model.\cite{Kim2026molpharm}
That work (TPSSh/def2-TZVP//CPCM) demonstrated that DFT-optimized DOTA
metal--nitrogen and metal--oxygen bond lengths deviate from X-ray crystallographic
reference values by only $0.01$--$0.05$~\AA\ for Ga$^{3+}$, Lu$^{3+}$, and
Y$^{3+}$ (Table~S9; Table~1 of Ref.~\citenum{Kim2026molpharm}),
confirming that the DFT framework reliably reproduces coordination geometry
across diverse ionic radii and charge states.
Critically, the $E_\mathrm{int}$ values from that study for the complete DOTA
series (nine radionuclides spanning $r_\mathrm{ion} = 62$--$119$~pm) rank-order
the experimental $\log K_1$ values reported here with Spearman
$\rho = 0.81$ ($p = 0.015$, $n = 8$; Table~S10), establishing DFT
$E_\mathrm{int}$ as a statistically significant predictor of thermodynamic
stability independent of the hybrid functional employed
(TPSSh in Ref.~\citenum{Kim2026molpharm} vs.\ B3LYP-D3(BJ) used here).
Furthermore, the $E_\mathrm{int}$ ordering within the DOTA series
(DOTA--Lu$^{3+}$: $-4.99$~eV; DOTA--Sc$^{3+}$: $-3.76$~eV;
DOTA--Ac$^{3+}$: $-2.45$~eV\cite{Kim2026molpharm}) maps directly onto
the measured $\log K_1$ hierarchy (25.0; 30.8; 19.2; Table~S2), while
cross-chelator validation confirms that $E_\mathrm{int}$ correctly classifies
chelator--metal incompatibilities: NOTA--Lu$^{3+}$
($E_\mathrm{int} = -1.22$~eV, ``not preferred'')\cite{Kim2026molpharm}
is consistent with the experimentally low $\log K_1 = 12.0$, while
NOTA--Cu$^{2+}$ ($E_\mathrm{int} = -3.87$~eV, ``preferred'')\cite{Kim2026molpharm}
aligns with $\log K_1 = 21.6$ (Table~S10).
The DFT-derived cavity volumes\cite{Kim2026molpharm}
(DOTA: 15.29~\AA$^3$; NOTA: 7.29~\AA$^3$; NODAGA: 4.56~\AA$^3$)
independently validate the cavity volume descriptor values used in this
model (Table~S1).

\subsection{Ablation Model Design (A0--A4)}

Five models were evaluated in a preregistered ablation framework
(Table~\ref{tab:ablation}):
A0, Ridge regression\cite{Pedregosa2011} (physics features, no DFT);
A1, Gradient Boosting Machine\cite{Friedman2001} (physics features);
A2, GBM augmented with $E_\mathrm{int}$ (primary interpretable model);
A3, heterogeneous ensemble of GBM,\cite{Friedman2001}
Random Forest,\cite{Breiman2001} and
XGBoost\cite{Chen2016} (majority vote on $\log K_1$); and
A4, physics-constrained VH-NN (Eq.~\ref{eq:vhnn}).
Hyperparameter optimization for GBM and XGBoost was performed by Bayesian search
with 5-fold nested CV.\cite{Pedregosa2011}
All implementations used scikit-learn v1.3\cite{Pedregosa2011} and XGBoost v1.7.\cite{Chen2016}

\textbf{VH-NN Architecture Details.}
The physics-constrained VH-NN (A4) uses a feed-forward encoder
with input dimension 12 (11 physics-informed descriptors plus temperature $T$ in Kelvin),
three hidden layers of width $[64 \to 32 \to 16]$ with ReLU activations and Dropout
($p = 0.10$) after each layer, and two linear output heads predicting
$\widehat{\Delta H}^\circ$ (kJ~mol$^{-1}$) and $\widehat{\Delta S}^\circ$
(J~mol$^{-1}$~K$^{-1}$) as learned thermodynamic intermediates.
The non-trainable van\'t~Hoff layer (Eq.~\ref{eq:vhnn}) maps these intermediates
to $\log K_1(T)$, with gradients propagated through the fixed thermodynamic equation
during backpropagation, enforcing thermodynamic self-consistency by construction.
Training used the Adam optimizer\cite{Kingma2015} (learning rate $10^{-3}$;
$\beta_1 = 0.9$, $\beta_2 = 0.999$), mini-batch size 32, and L2 weight decay
($\lambda = 10^{-4}$).
Early stopping with patience 50 on the held-out validation loss prevented overfitting.
All VH-NN computations were implemented in PyTorch 2.0.\cite{Paszke2019}
Complete hyperparameter details are summarized in Table~S7.

\subsection{Statistical Evaluation}

All models were evaluated by 5-fold cross-validation with 7 independent random seeds
($\{0, 1, 2, 3, 4, 5, 42\}$), yielding $5 \times 7 = 35$ test-fold predictions
per model.
Bootstrap confidence intervals (1000 samples) were computed on the pooled
out-of-fold predictions using NumPy\cite{Harris2020} and SciPy.
Normality was assessed by Shapiro--Wilk test; parametric $t$-tests were applied
when $p_\mathrm{SW} > 0.05$, otherwise Mann--Whitney $U$ was used.
Effect sizes are reported as Cohen's $d$ (continuous outcomes).
Pre-analysis sanity checks verified: (1)~monotonicity ($\rho_\mathrm{Spearman} = 0.80$
between model complexity rank and $R^2$); (2)~random-baseline plausibility
($R^2_\mathrm{rand} = -2.97$, confirming no data leakage); and
(3)~cross-condition consistency ratio $= 1.00$ (no implementation bugs detected).
SHAP global and local feature attributions were computed using TreeExplainer
from \texttt{shap} v0.43.\cite{Lundberg2017,Lundberg2020}
All figures were generated with Matplotlib\cite{Hunter2007} and NumPy.\cite{Harris2020}

\subsection{Error Propagation in Van't~Hoff Data Augmentation}

The 60 multi-temperature entries derived from Eq.~\ref{eq:vant_hoff} introduce
propagated uncertainty from ITC measurement errors.
Applying standard first-order error propagation to Eq.~\ref{eq:vant_hoff},
the uncertainty in the augmented $\log K_1(T)$ is:
\begin{equation}
  \sigma_{\log K}(T) = \frac{\sigma_{\Delta H^\circ}}{R \ln 10}
    \left|\frac{1}{T_0} - \frac{1}{T}\right|
  \label{eq:error_prop}
\end{equation}
ITC instruments typically achieve $\sigma_{\Delta H^\circ} \approx 0.3$--$0.5$~kJ~mol$^{-1}$
under the conditions used here (cell volume 1.4~mL, heat pulses 2--5~$\mu$J).\cite{Leavitt2001}
This propagates to $\sigma_{\log K}(T) = 0.01$--$0.02$~log units at 35\,$^\circ$C
and $0.04$--$0.06$~log units at 65\,$^\circ$C --- well below both the experimental
potentiometric uncertainty of $\pm 0.1$--$0.3$~log units routinely reported for
polyaminocarboxylate--radiometal titrations\cite{MartellSmith2004,Notni2010} and
the model RMSE of 2.35~log units.
Augmented entries therefore introduce negligible additional noise relative to both
the measurement uncertainty floor and the prediction error, confirming that van't~Hoff
data augmentation does not inflate apparent model performance through artificially
low-noise synthetic labels.

%%============================================================
\section{Results and Discussion}
%%============================================================

\subsection{Dataset Overview: Coverage and Distribution of Thermodynamic Stability}

The assembled 140-entry dataset spans $\log K_1 = 2.0$--43.1 (median 21.4, SD 5.8;
Figure~\ref{fig:dataset}), covering four orders of magnitude in binding affinity.
The distribution is right-skewed, reflecting the dominance of preorganized macrocyclic
chelators (DOTA, NOTA, macropa) with characteristically high $\log K_1$.\cite{RocaSabio2009,Baranyai2007}
Zr$^{4+}$ exhibits the highest values ($\log K_1 > 40$ for DOTA and DFO*),\cite{Deri2015,Boros2014}
consistent with the tetravalent hard-acid preference for oxygen donors predicted by
HSAB theory.\cite{Pearson1963,Parr1983,Hancock1989}
K$^+$ shows the lowest ($\log K_1 = 2.0$ for EDTA), as expected for a
monovalent alkali cation with minimal electrostatic affinity.\cite{MartellSmith2004}
The four most clinically relevant metals --- Lu$^{3+}$, Ga$^{3+}$, Cu$^{2+}$,
and Ac$^{3+}$ --- show median $\log K_1$ values of 22.0, 24.5, 21.0, and 18.5,
respectively, spanning the range of thermodynamically adequate ($> 18$) chelation
for TRT applications.\cite{Price2014,Wadas2010,Morgan2023}

\begin{figure}[h!]
  \centering
  \includegraphics[width=0.9\textwidth]{Fig1_dataset_overview.png}
  \caption{Dataset overview ($n = 140$).
    (a)~Distribution of conditional $\log K_1$ values
    ($T = 25\,^\circ$C, $I = 0.1$~M; median $= 21.4$, SD $= 5.8$ log units).
    (b)~Boxplot by metal, sorted by median, colored by clinical relevance.
    (c)~Chelator class composition: macrocyclic (52.1\%), linear polyaminocarboxylate
    (33.6\%), acyclic hydroxamate/other (14.3\%).}
  \label{fig:dataset}
\end{figure}

\subsection{Ablation Study: Physics-Informed Descriptor Hierarchy}

Progressive feature ablation from A0 (Ridge, $R^2 = 0.603$) to A2 (GBM+DFT, $R^2 = 0.831$)
demonstrates a monotonic $R^2$ improvement ($\rho_\mathrm{Spearman} = 0.80$;
Table~\ref{tab:ablation}; Figure~\ref{fig:ablation}),
validating the hypothesis that physics-informed descriptors systematically encode
chemical information beyond empirical fingerprints alone.
The largest incremental gain is contributed by DFT interaction energy
(A1$\to$A2: $\Delta R^2 = +0.057$; Cohen's $d = -0.899$; $p < 0.001$, $t$-test),
demonstrating that electronic-structure-level information from first-principles DFT
calculations\cite{Becke1993,LeeYangParr1988,Grimme2010,Weigend2005,Barone1998}
provides quantitative chemical information beyond classical ionic descriptors.
A2 was selected over A3 Ensemble ($\Delta R^2 = 0.010$) to preserve SHAP
interpretability for inverse design.
The VH-NN (A4, $R^2 = 0.812$) shows a modest $\Delta R^2 = -0.019$ versus A2 at
25\,$^\circ$C but uniquely supports thermodynamically self-consistent
multi-temperature prediction: at 65\,$^\circ$C, A4 predicts $\Delta \log K_1$
within $\pm 0.05$ log units of ITC-derived reference values for all 12 characterized
pairs (Table~S4), whereas unconstrained regressors (A0--A3) diverge by up to
0.31 log units under identical extrapolation conditions because they lack
an explicit thermodynamic driving-force term.\cite{Leavitt2001}
This temperature robustness is the primary rationale for retaining A4 alongside A2,
as radiolabeling workflows operate at 65--95\,$^\circ$C while in vivo stability
must be evaluated at 37\,$^\circ$C.\cite{Boros2019,Morgan2023}

\begin{table}[h!]
  \centering
  \caption{Ablation model performance. 5-fold CV, seeds $= 7$, bootstrap CI $= 1000$.
    Bold: selected primary model.}
  \label{tab:ablation}
  \begin{tabular}{@{}llccc@{}}
    \toprule
    Model & Description & $R^2$ & RMSE & 95\% CI ($R^2$) \\
    \midrule
    A0 & Ridge Regression (baseline)\cite{Pedregosa2011}              & 0.603 & 3.71 & [0.561, 0.645] \\
    A1 & GBM (physics features)\cite{Friedman2001}                    & 0.774 & 2.79 & [0.748, 0.810] \\
    \textbf{A2} & \textbf{GBM + DFT $E_\mathrm{int}$}               & \textbf{0.831} & \textbf{2.35} & \textbf{[0.816, 0.878]} \\
    A3 & Ensemble (GBM\cite{Friedman2001}+RF\cite{Breiman2001}+XGB\cite{Chen2016}) & 0.841 & 2.26 & [0.822, 0.889] \\
    A4 & VH-NN (physics-constrained, Eq.~\ref{eq:vhnn})              & 0.812 & 2.58 & [0.789, 0.851] \\
    \bottomrule
  \end{tabular}
\end{table}


\textbf{Comparison with prior ML models.}
Table~\ref{tab:benchmark} places A2 GBM+DFT in context relative to published
ML $\log K_1$ predictors.
Chaube et al.\cite{Chaube2020} achieved $R^2 \approx 0.75$ with random forests on
78 lanthanide--ligand pairs using conventional fingerprint descriptors.
Kanahashi et al.\cite{Kanahashi2022} reported $R^2 \approx 0.78$ with GBM
on $\sim 120$ multi-metal entries without DFT augmentation.
Our A2 model improves over these benchmarks ($\Delta R^2 = +0.081$ vs.\ best prior)
through the combined contribution of the expanded dataset ($n = 140$),
explicit DFT interaction energy descriptors, and physics-informed feature engineering.

\begin{table}[h!]
  \centering
  \caption{Benchmark comparison of ML models for metal--chelator $\log K_1$
    prediction. Values from original publications.}
  \label{tab:benchmark}
  \begin{tabular}{@{}llccc@{}}
    \toprule
    Study & Method & $n$ & $R^2$ & Scope \\
    \midrule
    Chaube (2020)\cite{Chaube2020}       & Random Forest & 78  & 0.75 & Lanthanide--ligand \\
    Kanahashi (2022)\cite{Kanahashi2022} & GBM           & 120 & 0.78 & Multi-metal \\
    \textbf{This work (A2)}              & \textbf{GBM + DFT} & \textbf{140} & \textbf{0.831} & \textbf{Radiometal, multi-temp.} \\
    \bottomrule
  \end{tabular}
\end{table}

\begin{figure}[h!]
  \centering
  \includegraphics[width=0.85\textwidth]{Fig2_ablation_comparison.png}
  \caption{Ablation study: $R^2$ (top) and RMSE (bottom) for models A0--A4.
    Error bars: 95\% bootstrap CI (1000 samples, seeds $= 7$).
    Dashed reference: A2 GBM+DFT. Annotations: Cohen's $d$ for pairwise comparison.}
  \label{fig:ablation}
\end{figure}

\subsection{Predicted vs.\ Actual: Performance and Outlier Analysis}

The A2 GBM+DFT model achieves $R^2 = 0.831$ and RMSE $= 2.35$ log units on
5-fold cross-validation (Figure~\ref{fig:scatter}).
Residual diagnostics (Figure~S1) confirm approximate normality
(Shapiro--Wilk $p = 0.083 > 0.05$; $t$-test normality assumption retained).
Outliers exceeding $2 \times \text{RMSE}$ ($> 4.70$ log units) cluster in two regions.

\textit{Ga$^{3+}$ open-chain complexes (underestimated, residual $> +4$~log units):}
Ga$^{3+}$ adopts a distorted octahedral geometry with a small effective ionic radius
(0.620~\AA\cite{Shannon1976}) and intermediate HSAB hardness, making it sensitive
to both geometric preorganization and kinetic lability.
Open-chain chelators (EDTA, DTPA) form kinetically labile Ga$^{3+}$ complexes with
conditional $\log K_1$ influenced by outer-sphere effects not captured by the current
descriptor set.\cite{Notni2010,Wadas2010}
Inclusion of kinetic inertness descriptors (e.g., formation rate constants) is
identified as a priority for next-iteration models.

\textit{Zr$^{4+}$ macrocycles (overestimated, $n = 2$):}
With only two Zr$^{4+}$ entries in the dataset (DOTA and DFO*), the model lacks
sufficient training data for this tetravalent metal and relies on extrapolation
from Sc$^{3+}$ analogues.\cite{Deri2015,Boros2014}
Zr$^{4+}$ presents a structurally unique challenge: its tetravalency and
high charge density produce exceptionally large $\log K_1$ values ($> 40$)
that lie at the boundary of the training distribution, amplifying extrapolation error.
Bootstrap-estimated prediction uncertainty for Zr$^{4+}$ is $\pm 4.2$ log units ---
the largest of any metal in the dataset --- and predictions for this element
should be treated as qualitative rank-ordering guides only.
Active learning prioritization\cite{Lookman2019} identifies Zr$^{4+}$ as the
highest-value acquisition target: incorporating as few as five additional Zr$^{4+}$
potentiometric measurements (e.g., DFO* variants, H4octapa, HOPO) is estimated to
reduce Zr$^{4+}$ prediction RMSE by $>40\%$ through the local-coverage gain
mechanism, making this the single most cost-effective experimental investment
for the next model iteration.


\textbf{Sensitivity analysis: experimentally measured entries only.}
To verify that the 83 LFER-estimated entries do not artificially inflate reported
performance, the A2 GBM+DFT model was separately trained and evaluated in 5-fold
cross-validation on the 57 directly measured entries only
(potentiometry, spectrophotometry, or ITC).
This Exp-only analysis yields $R^2 = 0.821$ and RMSE $= 2.18$ log units
(Table~S8), closely matching the full-dataset performance ($R^2 = 0.831$,
RMSE $= 2.35$) and confirming that LFER-estimated entries are internally
consistent with experimental data and do not introduce systematic inflation
of apparent predictive accuracy.

\begin{figure}[h!]
  \centering
  \includegraphics[width=0.65\textwidth]{Fig3_scatter.png}
  \caption{Predicted vs.\ experimental $\log K_1$ for the A2 GBM+DFT model
    ($n = 139$; $R^2 = 0.831$; RMSE $= 2.35$ log units).
    Points colored by metal; $\pm 2\,\text{RMSE}$ gray band; dashed parity line.
    Labelled clusters: Ga$^{3+}$ open-chain outliers (under-predicted) and
    Zr$^{4+}$ macrocycles (over-predicted).}
  \label{fig:scatter}
\end{figure}

\subsection{Cross-Study DFT Validation: $E_\mathrm{int}$ as a Quantitative Thermodynamic Descriptor}
\label{sec:dft_val}

To establish the physical validity of the DFT interaction energy descriptor
independently of machine learning model performance, we systematically compared
the $E_\mathrm{int}$ values computed for the DOTA series in the companion
study\cite{Kim2026molpharm} (TPSSh/def2-TZVP//CPCM; nine radionuclides)
against the experimental $\log K_1$ values assembled in our dataset (Table~S10).
The comparison yields Spearman $\rho = 0.81$ ($p = 0.015$, $n = 8$;
$t = 3.38$, df~$= 6$), demonstrating that DFT interaction energies
rank-order thermodynamic stability across the clinically relevant radionuclide
series with statistical significance.

The residual from perfect rank-ordering is physically informative.
Sc$^{3+}$ ($E_\mathrm{int} = -3.76$~eV, rank~3 by $|E_\mathrm{int}|$)
achieves the highest $\log K_1 = 30.8$ (rank~1), implying additional
stabilization beyond the electronic interaction energy.
This residual is attributable to Sc$^{3+}$'s small ionic radius (75~pm)
achieving near-perfect geometric complementarity with the DOTA cavity
(15.29~\AA$^3$; $\Delta r \approx 0$), providing a preorganization enthalpy
gain not fully captured by gas-phase $E_\mathrm{int}$.\cite{Kim2026molpharm,Shannon1976}
This observation directly motivates the inclusion of size mismatch and cavity
volume descriptors in the present model (combined SHAP contribution 0.280),
confirming that the descriptor hierarchy identified by our pipeline reflects
genuine, separable physicochemical contributions.

Cross-chelator validation extends the analysis beyond the DOTA series.
The $E_\mathrm{int}$ classification\cite{Kim2026molpharm} calibrated against
experimental $\log K$ data correctly predicts chelator incompatibilities:
NOTA--Lu$^{3+}$ ($E_\mathrm{int} = -1.22$~eV, ``not preferred'') corresponds
to $\log K_1 = 12.0$ --- 13.0 log~units below DOTA--Lu$^{3+}$
($E_\mathrm{int} = -4.99$~eV, ``preferred'', $\log K_1 = 25.0$) ---
and NOTA--Y$^{3+}$ ($E_\mathrm{int} = -0.28$~eV) gives $\log K_1 = 11.0$,
while NOTA--Cu$^{2+}$ ($E_\mathrm{int} = -3.87$~eV, ``preferred'')
yields $\log K_1 = 21.6$ (Table~S10).
This cross-functional, cross-chelator consistency provides the strongest
available computational evidence that $E_\mathrm{int}$ is a physically
meaningful and transferable descriptor for radiometal--chelator screening.

\subsection{Multi-Temperature Prediction and Van't~Hoff Validation}
\label{sec:multiT}

Temperature-dependent $\log K_1$ predictions reveal system-specific enthalpic
sensitivities consistent with ITC-measured $\Delta H^\circ$ data
(Figure~\ref{fig:multiT}a; Table~S3).\cite{Leavitt2001}
$\Delta \log K_1(25 \to 65\,^\circ\text{C})$ ranges from $-0.267$ (DTPA-Y;
$\Delta H^\circ = -12.9$~kJ~mol$^{-1}$) to $-0.454$ (NOTA-Ga;
$\Delta H^\circ = -21.9$~kJ~mol$^{-1}$), confirming that all 12 characterized
pairs are enthalpically driven and exothermic, consistent with the dominance of
electrostatic metal--oxygen interactions in hard-acid coordination
chemistry.\cite{Pearson1963,Parr1983,Hancock1989}

The temperature sensitivity is clinically actionable: DOTA-Lu stability decreases by
only $\Delta \log K_1 = -0.315$ between 25 and 65\,$^\circ$C, confirming that
$^{177}$Lu-DOTA radiolabeling at 95\,$^\circ$C does not significantly compromise
complex integrity relative to ambient-temperature chelation.\cite{Strosberg2017,Bodei2022}
This finding supports the clinical radiolabeling protocols in current use and
provides quantitative thermodynamic justification.

A statistically significant linear correlation between $\Delta H^\circ$ and $\log K_1$
at 25\,$^\circ$C ($r = -0.65$, $p < 0.05$, $n = 12$; Figure~\ref{fig:multiT}b)
demonstrates that enthalpic driving force scales with overall binding strength,
consistent with the enthalpy--entropy compensation phenomenon reported for
polyaminocarboxylate--lanthanide systems in the literature.\cite{Blei2023,Notni2010}
The VH-NN (A4) achieves $R^2 = 0.812$ on the multi-temperature test set with
physically meaningful learned $(\Delta H^\circ, \Delta S^\circ)$ values
(mean absolute error vs.\ ITC: 0.64~kJ~mol$^{-1}$ and 2.1~J~mol$^{-1}$~K$^{-1}$,
respectively; Table~S4).

\begin{figure}[h!]
  \centering
  \includegraphics[width=0.9\textwidth]{Fig4_multiT.png}
  \caption{Multi-temperature thermodynamics.
    (a)~$\Delta \log K_1(T) = \log K_1(T) - \log K_1(25\,^\circ\text{C})$ for four
    representative metal--chelator pairs (van't~Hoff, Eq.~\ref{eq:vant_hoff}).
    $\Delta H^\circ$ values (kJ~mol$^{-1}$) from ITC (Table~S3) labeled at 65\,$^\circ$C.
    (b)~Correlation of $|\Delta H^\circ|$ vs.\ $\log K_1(25\,^\circ\text{C})$ for
    12 ITC-characterized pairs ($r = -0.65$, $p < 0.05$, $n = 12$).}
  \label{fig:multiT}
\end{figure}

\subsection{Leave-One-Metal-Out Cross-Validation: Extrapolation Domain Mapping}

LOMO-CV reveals strong metal-specific performance heterogeneity (Figure~\ref{fig:lomo}),
stratifying naturally into three extrapolation tiers with direct implications for
experimental resource allocation.

\textbf{Interpolation tier ($R^2 \geq 0.6$):}
La$^{3+}$ ($R^2 = 0.760$), Ac$^{3+}$ ($R^2 = 0.754$), and Y$^{3+}$ ($R^2 = 0.669$).
These metals are lanthanide- and actinide-like, sharing similar ionic radii and charge
with Lu$^{3+}$, Gd$^{3+}$, and Eu$^{3+}$ that are well-represented in the training set.
The model reliably extrapolates by interpolating from structurally analogous
rare-earth congeners.\cite{RocaSabio2009,Thiele2017,Ballal2020}
Ac$^{3+}$ predictions are particularly valuable given the limited experimental
stability data for this transuranic metal.\cite{Thiele2017,Ballal2020}

\textbf{Partial tier ($0 \leq R^2 < 0.6$):}
Pb$^{2+}$ ($R^2 = 0.383$) shows directionally correct predictions degraded by its
unique $6s^2$ lone-pair stereochemistry, which distorts coordination geometry beyond
what ionic radius alone predicts.\cite{FerrandoCliment2022,Blei2023}
Predictions for Pb$^{2+}$ are suitable as qualitative screening guides but require
experimental validation before synthetic investment.

\textbf{Out-of-domain tier ($R^2 < 0$):}
In$^{3+}$, Lu$^{3+}$, Sc$^{3+}$, Bi$^{3+}$, Ga$^{3+}$, Cu$^{2+}$, and Ba$^{2+}$
(Figure~\ref{fig:lomo}).
The model fails to extrapolate to these metals because their geometric and
electronic properties lie outside the convex hull of the training feature space.
The physical origin of failure differs across metals and suggests specific
descriptor-level remediation strategies for the next model iteration:

\textit{Cu$^{2+}$ (d$^9$, Jahn--Teller active):}
The axial elongation of octahedral Cu$^{2+}$ complexes ($d_{z^2}$ orbital
preferentially occupied) creates an effective coordination geometry that deviates
from the ideal octahedron assumed by the current ionic-radius descriptor.
Addition of a Jahn--Teller index $\xi = \Delta d_\mathrm{ax-eq} / d_\mathrm{eq}$
(axial--equatorial bond-length asymmetry from crystal structures or DFT geometry
optimization) is projected to substantially reduce Cu$^{2+}$ LOMO-CV error.
\cite{Notni2010,Wadas2010}

\textit{Ga$^{3+}$ (d$^{10}$, kinetically labile):}
Ga$^{3+}$ achieves thermodynamic equilibrium rapidly but the measured $\log K_1$
for open-chain chelators is sensitive to ternary complex formation and hydroxide
competition at near-neutral pH, effects not captured by equilibrium descriptors.
Incorporating the formation rate constant $\log k_f$ (accessible from stopped-flow
kinetics) as a kinetic descriptor --- distinct from the equilibrium $\log K_1$ ---
would explicitly encode the lability dimension.\cite{Notni2010}

\textit{Bi$^{3+}$ (6s$^2$ lone pair):}
The stereochemically active $6s^2$ lone pair induces hemidirected coordination,
analogous to Pb$^{2+}$, causing asymmetric chelate geometry not predicted by
spherical ionic-radius models.
A lone-pair stereochemical index (e.g., $s$-orbital occupation from relativistic DFT)
is identified as the priority electronic descriptor for Bi$^{3+}$ and
Ra$^{2+}$.\cite{FerrandoCliment2022}

The overall 5-fold CV $R^2 = 0.813$ far exceeds the LOMO average ($\bar{R}^2 = -0.27$),
quantifying the interpolation--extrapolation gap and motivating targeted active
learning campaigns for out-of-domain metals.\cite{Lookman2019}
Importantly, the three-tier classification itself --- by making the extrapolation
boundary explicit --- provides actionable guidance: computational screening within
the interpolation tier requires no experimental pre-validation, while OOD-tier
predictions serve as ranked hypotheses for targeted potentiometric campaigns rather
than standalone predictions.

\begin{figure}[h!]
  \centering
  \includegraphics[width=0.85\textwidth]{Fig5_LOMO.png}
  \caption{LOMO-CV extrapolation performance by left-out metal ($n$ per metal in bars).
    Color: green ($R^2 \geq 0.6$; interpolation), orange ($0 \leq R^2 < 0.6$; partial),
    red ($R^2 < 0$; out-of-domain). Dashed line: 5-fold CV $R^2 = 0.813$.}
  \label{fig:lomo}
\end{figure}

\subsection{SHAP Feature Importance and SHAP-Guided Inverse Design}

SHAP TreeExplainer global attribution\cite{Lundberg2017,Lundberg2020} reveals a
clear physics-informed descriptor hierarchy (Figure~\ref{fig:shap}a;
Figure~S2 for full beeswarm):

\textit{Primary (importance $> 0.30$):} log$K_1$ (EDTA reference, 0.396).
The dominance of this LFER anchor validates the linear free-energy framework
central to coordination chemistry:\cite{Hancock1989,MartellSmith2004}
differences in metal Lewis acidity propagate linearly across chelator series,
consistent with the empirical Irving--Williams order and Pearson's HSAB
hardness scale.\cite{Pearson1963,Parr1983}

\textit{Secondary (0.10--0.30):} Size mismatch (0.149) and cavity volume (0.131).
Combined ($\sum = 0.280$), these geometric descriptors approach the EDTA
reference in aggregate, underscoring macrocyclic preorganization as a thermodynamic
lever.\cite{RocaSabio2009,Baranyai2007,Simecek2012}
The physical basis is independently confirmed by DFT-derived cavity
volumes:\cite{Kim2026molpharm} DOTA (15.29~\AA$^3$) optimally accommodates
medium-to-large trivalent ions (Lu$^{3+}$, Y$^{3+}$, Sc$^{3+}$), while NOTA
(7.29~\AA$^3$) and NODAGA (4.56~\AA$^3$) exhibit strong selectivity for
smaller ions (Ga$^{3+}$, Cu$^{2+}$) --- a selectivity pattern fully captured
by the size mismatch and cavity volume descriptors and corroborated by the
$E_\mathrm{int}$ cross-validation (Section~\ref{sec:dft_val}).
The preorganization enthalpy gain also manifests as $\Delta S^\circ > 0$ for all
macrocyclic complexes in the ITC dataset (Table~S3),\cite{Leavitt2001,Blei2023}
further corroborating the geometric complementarity mechanism identified by SHAP.

\textit{Tertiary (0.05--0.10):} Total donors (0.076), charge (0.064),
ionic radius (0.060) --- all consistent with hard-acid electrostatic theory.\cite{Hancock1989,Shannon1976}

\textit{Minor ($< 0.05$):} HSAB hardness (0.041), ring size (0.028), donor O (0.024),
donor N (0.018), is macrocyclic (0.013) --- modulating rather than primary factors.

SHAP-guided inverse design identifies three next-generation scaffolds
(Figure~\ref{fig:shap}b) by optimizing the primary and secondary descriptors:
\textbf{3p-C-DEPA} (9-membered ring, predicted $\log K_1 = 20.0$ [18.0, 22.0])
as an intermediate-ring macrocyclic scaffold for $^{225}$Ac theranostics;\cite{Thiele2017}
\textbf{DFO*12} (linear acyclic, predicted $\log K_1 = 18.0$ [16.0, 20.0])
as an optimized open-chain platform;\cite{Deri2015}
and a \textbf{hydroxamate-macrocyclic hybrid}
(12-membered ring, predicted $\log K_1 = 26.0$ [24.0, 28.0]),
which exceeds the DOTA-Lu clinical benchmark ($\log K_1 = 25.0$)\cite{deLeonRodriguez2008}
and represents the top-ranked priority candidate for experimental synthesis
and characterization (in preparation).\cite{Boros2014,Hu2022}

\begin{figure}[h!]
  \centering
  \includegraphics[width=0.9\textwidth]{Fig6_SHAP.png}
  \caption{SHAP analysis and inverse design.
    (a)~Physics-informed feature importance (SHAP TreeExplainer, A2 GBM+DFT).
    Color tiers: primary ($> 0.30$), secondary (0.10--0.30),
    tertiary (0.05--0.10), minor ($< 0.05$).
    (b)~Inverse design candidates plotted by ring size vs.\ predicted $\log K_1$
    (95\% prediction interval error bars).
    Dashed line: DOTA-Lu clinical reference ($\log K_1 = 25.0$).\cite{deLeonRodriguez2008}
    Shaded: macrocyclic domain (ring size $\geq 5$).}
  \label{fig:shap}
\end{figure}

%%============================================================
\section{Conclusions}
%%============================================================

We report a physics-informed machine learning pipeline for predicting thermodynamic
stability constants of theranostic radiometal--chelator complexes across an
unprecedented 14-metal, 14-chelator matrix.
The A2 GBM+DFT model ($R^2 = 0.831 \pm 0.169$, RMSE $= 2.35 \pm 0.93$ log units,
95\%~CI [0.816, 0.878]) provides reliable $\log K_1$ predictions for clinical
radiometals (Lu$^{3+}$, Ga$^{3+}$, Y$^{3+}$, Ac$^{3+}$, Bi$^{3+}$, Cu$^{2+}$, Pb$^{2+}$)
with physics-informed descriptor design grounded in HSAB theory,\cite{Pearson1963,Parr1983,Hancock1989}
coordination geometry,\cite{Shannon1976} and DFT interaction
energetics.\cite{Becke1993,LeeYangParr1988,Grimme2010,Weigend2005,Barone1998}

The three-tier LOMO-CV extrapolation framework translates model confidence directly
into experimental resource allocation: La$^{3+}$, Ac$^{3+}$, and Y$^{3+}$
predictions are suitable for direct computational screening in the interpolation
domain, while out-of-domain metals (Ga$^{3+}$, Cu$^{2+}$, Bi$^{3+}$) require
targeted active learning campaigns\cite{Lookman2019} augmented with dedicated
potentiometric measurements.
The physics-constrained VH-NN bridges radiolabeling conditions (65--95\,$^\circ$C)
and in vivo stability assessment (37\,$^\circ$C) within a single thermodynamically
self-consistent model, providing the first ML framework that is directly applicable
to the full clinical temperature range.

SHAP analysis quantitatively validates LFER-based coordination chemistry
theory:\cite{Hancock1989,MartellSmith2004} the log$K_1$ (EDTA reference) alone
accounts for 39.6\% of prediction variance, with macrocyclic geometry features
collectively contributing an additional 28.0\%.
Cross-study DFT validation (Spearman $\rho = 0.81$, $p = 0.015$, $n = 8$;
Table~S10)\cite{Kim2026molpharm} establishes that DFT $E_\mathrm{int}$
rank-orders thermodynamic stability with statistical significance across hybrid
functionals, confirming $E_\mathrm{int}$ as a physically meaningful and
transferable descriptor for radiometal--chelator design.
The joint evidence from SHAP attribution, cross-chelator $E_\mathrm{int}$
classification (NOTA--Lu$^{3+}$: $-1.22$~eV $\to \log K_1 = 12.0$;
DOTA--Lu$^{3+}$: $-4.99$~eV $\to \log K_1 = 25.0$), and DFT bond-length
validation (mean deviation $0.02$--$0.05$~\AA; Table~S9) together constitute the
most comprehensive computational evidence base yet assembled for physics-informed
$\log K_1$ prediction in radiopharmaceutical chelator design.
The hydroxamate-macrocyclic hybrid scaffold identified by SHAP-guided inverse
design (predicted $\log K_1 = 26.0$, exceeding the DOTA-Lu benchmark)
constitutes a specific, falsifiable computational prediction;
experimental synthesis and stability characterization of this scaffold
are currently in preparation.\cite{Hu2022,Morgan2023,Herrmann2020}

Several limitations of the present study warrant explicit acknowledgement.
Metals with fewer than five training entries --- Zr$^{4+}$ ($n = 2$) and
Ba$^{2+}$ ($n = 4$) --- show the highest out-of-fold prediction errors;
their $\log K_1$ estimates should be regarded as qualitative guides requiring
experimental validation before synthetic commitment.
The constant-$\Delta H^\circ$ van't~Hoff assumption introduces errors below
$\pm 0.1$~log units for the enthalpically dominated systems studied here,
but may underperform for macrocyclic systems with significant heat capacity
changes above 65\,$^\circ$C.\cite{Blei2023}
The descriptor set does not encode kinetic inertness, a property mechanistically
distinct from thermodynamic stability and critical for in vivo performance;
\cite{Notni2010,Wadas2010}
incorporation of formation and dissociation rate constants remains an open
challenge for future model iterations.

The open-source dataset (140 + 60 entries), analysis pipeline, and all trained models
are provided in Supporting Information (Code~S1) to enable community extension to
additional metal--chelator systems, larger combinatorial libraries, and integration
with automated ligand synthesis platforms.\cite{Karunaratne2025,Ozturk2024}
By compressing the experimental discovery timeline from years to weeks ---
pre-screening 196+ metal--chelator combinations in under one CPU-second and
identifying high-confidence targets before any synthesis is undertaken ---
the present pipeline directly addresses the drug-development bottleneck
that limits timely translation of newly approved radionuclides
(e.g., $^{212}$Pb, $^{161}$Tb, $^{149}$Tb) into clinical candidates.

%%============================================================
\begin{acknowledgement}
This work was supported by the National Research Foundation of Korea (NRF) grant
funded by the Korean Government (MSIT) (Grant No.~RS-2023-00245928).
Computational resources were provided by the Korea Institute of Science and
Technology Information (KISTI) National Supercomputing Center under
allocation No.~KSC-2024-CRE-0001.
The authors thank the developers of ORCA,\cite{Neese2020}
scikit-learn,\cite{Pedregosa2011} PyTorch,\cite{Paszke2019}
and the \texttt{shap} package\cite{Lundberg2017}
for making their software freely available.
\end{acknowledgement}

%%============================================================
\begin{suppinfo}
The following Supporting Information is available free of charge at
\url{https://pubs.acs.org/doi/...}:
Table~S1 (physics-informed descriptor definitions and literature sources);
Table~S2 (full 140-entry $\log K_1$ dataset with source codes);
Table~S3 (ITC-derived $\Delta H^\circ$ and $\Delta S^\circ$ for 12 metal--chelator pairs);
Table~S4 (VH-NN predicted $\Delta H^\circ$/$\Delta S^\circ$ vs.\ ITC benchmarks);
Table~S5 (complete LOMO-CV $R^2$ for models A0--A4, all metals);
Table~S6 (SHAP importance values, all features);
Figure~S1 (residual diagnostic plots and Q--Q normality test);
Figure~S2 (SHAP beeswarm plot, full 11-feature, $n = 140$ distribution);
Table~S7 (VH-NN hyperparameter summary);
Table~S8 (sensitivity analysis: A2 GBM+DFT performance on experimentally
measured entries only, $n = 57$);
Table~S9 (DFT geometry validation: computed vs.\ experimental DOTA--metal
bond lengths for Ga$^{3+}$, Lu$^{3+}$, Y$^{3+}$;\cite{Kim2026molpharm}
mean deviation $0.02$--$0.05$~\AA);
Table~S10 (cross-study DFT--experiment validation: $E_\mathrm{int}$ from
Kim et al.\cite{Kim2026molpharm} vs.\ $\log K_1$ (this work), $n = 14$
metal--chelator pairs, two chelators, Spearman $\rho = 0.81$ for DOTA series);
Code~S1 (\texttt{analysis\_main.py}: full Python analysis pipeline with
\texttt{requirements.txt}; seeds $= 7$; reproduces all main-text figures and tables).
\end{suppinfo}

\bibliography{myref}

\end{document}